% ---------------------------------------------------------------------
%  CHƯƠNG 3 — KẾT QUẢ, ĐÁNH GIÁ VÀ KẾT LUẬN
% ---------------------------------------------------------------------
\chapter{Kết quả, đánh giá và kết luận}
\label{ch:ketqua}

\begin{muctieu}
Chương này đối chiếu mục tiêu với sản phẩm đã triển khai, trình bày bằng
chứng kiểm thử, thảo luận giới hạn và tổng kết hướng phát triển.
\end{muctieu}

\section{Kết quả đạt được}

Các mục tiêu ở Chương~\ref{ch:modau} được đối chiếu tại
Bảng~\ref{tab:doichieu}. Sản phẩm hiện bao gồm ứng dụng web, API và lược đồ
cơ sở dữ liệu; các luồng tạo câu hỏi, lập đề, duyệt đề, thi trực tuyến và
lớp học có mã nguồn cùng kiểm thử. Thành phần tạo câu hỏi hỗ trợ gọi mô
hình ngôn ngữ tùy cấu hình và có nhánh dự phòng. Các thành phần này chứng
minh mức độ hoàn thiện chức năng của nguyên mẫu, không chứng minh chất lượng
chuyên môn của câu hỏi do máy tạo.

\begin{table}[H]
\centering
\caption{Đối chiếu mục tiêu và kết quả triển khai.}
\label{tab:doichieu}
\begin{tabular*}{\textwidth}{@{\extracolsep{\fill}}p{0.31\textwidth}p{0.53\textwidth}p{0.10\textwidth}}
\toprule
\rowcolor{gray100}\bfseries\color{daunavy}Mục tiêu & \bfseries\color{daunavy}Bằng chứng trong dự án & \bfseries\color{daumaroon}Trạng thái \\
\midrule
Ứng dụng web nhiều thành phần & React/TypeScript, Express API, PostgreSQL và Docker Compose. & Đạt \\
Hỗ trợ sinh câu hỏi từ học liệu & Trích xuất PDF/DOCX, tích hợp OpenAI tùy chọn, sinh dự phòng và kiểm tra đầu ra. & Đạt một phần \\
Quản lý ngân hàng và đề & CRUD câu hỏi/đề, chọn theo loại và độ khó, tạo mã đề, quy trình duyệt. & Đạt \\
Tổ chức thi và lớp học & Mã phòng/lớp, lưu lượt làm bài, giám sát, tiếp tục bài và tự chấm trắc nghiệm. & Đạt \\
Kiểm thử chức năng và build & 35/35 bài kiểm thử API đạt; frontend build/type-check thành công. & Đạt \\
Đánh giá định lượng chất lượng AI & Chưa có tập chuẩn gán nhãn, phép đo chất lượng sinh hoặc benchmark OCR. & Chưa đạt \\
\bottomrule
\end{tabular*}
\end{table}

\section{Độ đo và bằng chứng kiểm thử}

Vì đồ án không huấn luyện mô hình và chưa có tập kiểm thử chuẩn cho câu hỏi
hay OCR, không thể báo cáo \en{accuracy}, precision, recall, F1 hoặc độ
chính xác OCR như kết quả thực nghiệm. Bằng chứng định lượng hiện có là số
bài kiểm thử phần mềm và kết quả build được trình bày trong
Bảng~\ref{tab:kiemthu}.

\begin{table}[H]
\centering
\caption{Kết quả xác minh phần mềm trong lần kiểm tra báo cáo.}
\label{tab:kiemthu}
\begin{tabular*}{\textwidth}{@{\extracolsep{\fill}}p{0.36\textwidth}p{0.36\textwidth}p{0.20\textwidth}}
\toprule
\rowcolor{gray100}\bfseries\color{daunavy}Hạng mục & \bfseries\color{daunavy}Kết quả & \bfseries\color{daumaroon}Đánh giá \\
\midrule
Kiểm thử API & 35 bài; 35 đạt, 0 lỗi, 0 bỏ qua. & Đạt \\
Build frontend & TypeScript và Vite build thành công. & Đạt \\
Sinh câu hỏi không có LLM & Có bài kiểm thử xác nhận nhánh sinh dự phòng trả ứng viên. & Đạt theo kiểm thử chức năng \\
Chất lượng nội dung AI & Chưa đối chiếu với bộ câu hỏi chuẩn hoặc đánh giá chuyên gia có thang điểm. & Chưa đo \\
Hiệu năng và tải & Chưa có benchmark thời gian phản hồi hoặc tải đồng thời. & Chưa đo \\
\bottomrule
\end{tabular*}
\end{table}

Backend trả về trường \code{confidence} theo nhánh xử lý, nhưng giá trị này
được gán trong mã nguồn và chưa được hiệu chuẩn bằng tập dữ liệu có nhãn.
Do đó, báo cáo không diễn giải trường này thành xác suất đúng hay chỉ số
đánh giá mô hình. Kiểm thử tự động tập trung vào cấu trúc đầu ra và nghiệp
vụ API; giảng viên vẫn cần kiểm tra nội dung câu hỏi, đáp án và bằng chứng
trước khi đưa vào đề thật.

\subsection{Ma trận kiểm thử chức năng}
Bảng~\ref{tab:ma-tran-kiem-thu} tóm tắt các nhóm hành vi được kiểm tra
trong bộ test hiện có. Mục đích là chỉ ra loại lỗi có thể được phát hiện,
không thay thế báo cáo từng lần chạy của CI. Tất cả 35 test hiện hữu đều
được thực thi thành công trong lần xác minh. Các yêu cầu không xuất hiện
trong test hiện hữu được giữ ở trạng thái chưa xác minh.
\begin{table}[H]
\centering
\caption{Ma trận kiểm thử chức năng theo rủi ro nghiệp vụ.}
\label{tab:ma-tran-kiem-thu}
\begin{tabularx}{\textwidth}{>{\bfseries\color{daunavy}}p{0.12\textwidth} X p{0.28\textwidth}}
\toprule
\rowcolor{gray100}Mã & Ca kiểm thử/điều kiện & Kết quả xác minh \\
\midrule
T01 & Health check và CORS preflight. & Đạt trong bộ test. \\
T02 & Cập nhật vai trò người dùng; token hiện hành và quyền được làm mới. & Đạt trong bộ test. \\
T03 & Sinh viên không truy cập phân tích quản trị; đăng ký công khai không tự cấp teacher/admin. & Đạt trong bộ test. \\
T04 & CRUD tài khoản quản trị; phân trang và bulk import. & Đạt trong bộ test. \\
T05 & Tạo, lưu và cập nhật đề; sinh mã đề với đáp án ánh xạ đúng. & Đạt trong bộ test. \\
T06 & Phân tích tài liệu có câu hỏi/đáp án; dữ liệu rỗng bị từ chối. & Đạt trong bộ test. \\
T07 & Không có LLM vẫn có ứng viên dự phòng; câu hỏi phải phù hợp văn bản nguồn. & Đạt theo ca hiện có. \\
T08 & Câu hỏi trắc nghiệm có lựa chọn riêng biệt, đáp án và căn cứ nguồn. & Đạt theo ca hiện có. \\
T09 & Phòng thi random đề, lưu bài, tự chấm trắc nghiệm và tiếp tục lượt. & Đạt trong bộ test. \\
T10 & Lớp học: tạo, tham gia bằng mã, danh sách thành viên và kiểm soát quyền. & Đạt trong bộ test. \\
T11 & OCR tiếng Việt, kiểm thử tải, tương phản/khả năng tiếp cận và trình duyệt thiết bị thật. & Chưa có bằng chứng. \\
\bottomrule
\end{tabularx}
\end{table}

\subsection{Diễn giải các loại kết quả}
Kiểm thử API kiểm chứng mã phản hồi, dữ liệu trả về và một phần trạng thái
PostgreSQL. Kiểm thử tích hợp chạm vào nhiều thành phần cùng lúc, nên có
thể phát hiện lỗi phối hợp mà test đơn vị của một hàm không phát hiện.
Tuy vậy, các bài test được viết cho những kịch bản xác định, không bao phủ
tất cả đầu vào, trình duyệt, nhà cung cấp OAuth hay cấu hình vận hành.

Lệnh build frontend xác nhận TypeScript không báo lỗi trong phạm vi
project build và bundler đóng gói được các tài nguyên. Nó không đo khả
năng sử dụng hoặc độ chính xác giao diện. Không có lượt test người dùng
được tiến hành, do vậy báo cáo không tự đưa ra kết luận về mức độ hài lòng
hoặc giảm thời gian biên soạn của giảng viên.

\subsection{Thiết kế thực nghiệm đánh giá AI đề xuất}
Để đánh giá chất lượng sinh câu hỏi trong tương lai, nên lập tập kiểm thử
đại diện cho nhiều môn và định dạng tài liệu, sau đó giảng viên độc lập
gán nhãn đúng/sai cho câu hỏi, đáp án, bằng chứng, độ rõ ràng, mức độ phù
hợp và độ khó. Một phần mẫu nên được hai người chấm để tính mức nhất quán
liên giám khảo; bất đồng cần có quy trình phân xử. Cần tách tài liệu đã
dùng để chỉnh prompt khỏi tập đánh giá cuối cùng nhằm hạn chế rò rỉ thông
tin vào phép đo.

Các câu hỏi trắc nghiệm có thể báo cáo tỷ lệ đáp án đúng được chuyên gia
xác nhận, tỷ lệ bằng chứng hỗ trợ phát biểu, tỷ lệ lựa chọn hợp lệ và
độ trùng lặp. Kết quả OCR có thể đo tỷ lệ lỗi ký tự và từ trên ảnh được
gán nhãn; điểm số OCR không thay cho đánh giá hậu quả của một lỗi OCR lên
câu hỏi sinh ra. Đánh giá phải phân tầng theo ngôn ngữ, loại tài liệu và
chất lượng scan, đồng thời công bố số mẫu và khoảng tin cậy.

\begin{table}[H]
\centering
\caption{Đề cương đo lường dành cho một đánh giá AI độc lập trong tương lai.}
\label{tab:ai-eval-plan}
\begin{tabularx}{\textwidth}{>{\bfseries\color{daunavy}}p{0.21\textwidth} X X}
\toprule
\rowcolor{gray100}Khía cạnh & Cách đo đề xuất & Đối tượng đánh giá \\
\midrule
Độ đúng nội dung & Tỷ lệ phát biểu/đáp án được chuyên gia xác nhận trên tập chưa dùng để phát triển. & Mỗi câu sinh và đáp án. \\
Bám nguồn & Tỷ lệ câu có bằng chứng nguyên văn, định vị đúng và đủ để hỗ trợ đáp án. & Cặp câu hỏi -- nguồn. \\
Chất lượng sư phạm & Thang điểm rõ ràng, một đáp án đúng, distractor hợp lý, ngôn ngữ và mức độ nhận thức. & Giảng viên chấm độc lập. \\
Độ đa dạng & Tỷ lệ trùng lặp nội dung và phân bố ý/chủ đề trong tập câu hỏi. & Tập đề sinh từ nhiều tài liệu. \\
OCR & Character/Word Error Rate trên ảnh quét có transcript chuẩn. & Từng nhóm ngôn ngữ và chất lượng ảnh. \\
Độ trễ/chi phí & Thời gian p50/p95, lỗi/time-out và chi phí theo tài liệu/câu hỏi. & Pipeline và từng nhà cung cấp. \\
\bottomrule
\end{tabularx}
\end{table}

Mỗi độ đo cần được báo cáo theo phương pháp đã định trước; không gộp nhiều
nhà cung cấp hoặc nhiều loại tài liệu thành một con số duy nhất nếu phân bố
mẫu khác nhau. Trường \code{confidence} hiện được gán theo nhánh xử lý và
chưa được hiệu chuẩn, nên không thể thay thế xác nhận chuyên gia hoặc
khoảng tin cậy thống kê.

\section{Bàn luận}

Điểm mạnh của thiết kế là có thể chạy quy trình nghiệp vụ liên tục từ tài
liệu và ngân hàng câu hỏi tới đề thi, lượt thi và kết quả; nhánh dự phòng
cũng cho phép kiểm thử luồng sinh mà không phụ thuộc khóa dịch vụ ngoài.
Việc lưu bằng chứng nguồn và hỗ trợ giảng viên rà soát là phù hợp với yêu
cầu truy xuất nội dung, nhưng không loại bỏ hoàn toàn nguy cơ câu hỏi sai
hoặc trích xuất sai.

Kết quả 35/35 bài đạt phản ánh phạm vi các ca kiểm thử đã viết. Nó không
đồng nghĩa hệ thống không còn lỗi, cũng không đo độ chính xác nội dung do
AI tạo. Build frontend thành công xác nhận khả năng biên dịch tại môi trường
đã kiểm tra, nhưng chưa xác nhận triển khai lên máy chủ thực tế.

\subsection{Đánh giá theo các vai trò}
Đối với giảng viên, việc gộp ngân hàng, bộ lọc ma trận và lịch sử đề làm
nền tảng cho khả năng tái sử dụng nội dung. Giá trị thực tiễn sẽ phụ thuộc
vào độ chính xác của căn cứ và thời gian phải chỉnh sửa câu AI; đồ án chưa
thực hiện đo trước/sau để kết luận tiết kiệm thời gian. Đối với sinh viên,
lưu vị trí và cho tiếp tục bài làm hỗ trợ tính liên tục, còn mã truy cập
tách việc dự thi khỏi việc gửi liên kết trực tiếp đến từng tài khoản.
Đối với quản trị viên, quản lý vai trò tập trung thuận tiện nhưng làm tăng
hậu quả nếu tài khoản có quyền cao bị chiếm dụng.

\subsection{Đánh đổi kiến trúc}
Tích hợp logic AI trong backend giúp sử dụng cơ sở dữ liệu và phân quyền
đã có, giảm số dịch vụ phải vận hành. Mặt trái là tác vụ trích xuất/gọi
LLM tốn thời gian có thể giữ request HTTP lâu và dùng bộ nhớ khi file được
đọc trong memory. Việc tách AI thành worker hoặc hàng đợi có thể cải thiện
độ bền và khả năng quan sát ở giai đoạn sau, nhưng đòi hỏi cơ chế trạng
thái job, retry an toàn, idempotency và giới hạn chi phí; chưa nên xem là
thành phần đã hiện thực.

\section{Hạn chế}

\begin{itemize}[leftmargin=1.6em, itemsep=3pt]
  \item Chưa có bộ dữ liệu chuẩn để đo độ đúng, độ bao phủ và độ đa dạng của
        câu hỏi sinh tự động.
  \item Chất lượng đầu ra phụ thuộc mô hình/API ngoài khi được cấu hình; bộ
        sinh dự phòng dựa trên quy tắc có thể tạo câu hỏi kém tự nhiên hoặc
        chưa đủ sâu.
  \item OCR chịu ảnh hưởng của chất lượng scan và bố cục; nhánh hiện tại
        dùng dữ liệu nhận dạng tiếng Anh, chưa có benchmark cho tài liệu
        tiếng Việt hoặc tập ảnh được gán nhãn.
  \item Câu tự luận cần được giảng viên chấm; giao diện có trường rubric và
        phản hồi nhưng báo cáo chưa xác nhận một bộ chấm tự động bằng AI.
  \item Chưa thực hiện kiểm thử tải, kiểm thử người dùng, đánh giá bảo mật
        triển khai và đo khả năng mở rộng.
  \item Các thư mục \code{ai-services/} và \code{ml-models/} mới là khung dự
        kiến; OCR và sinh câu hỏi hiện được gọi từ backend API.
\end{itemize}
